\documentclass[12pt]{article}

\usepackage[T1]{fontenc}
\usepackage[utf8]{inputenc}
\usepackage[brazil]{babel}
\usepackage{lmodern}
\usepackage{microtype}
\usepackage[a4paper,margin=2.5cm]{geometry}
\usepackage{setspace}
\usepackage{indentfirst}

\setstretch{1.15}

\title{Aspectos Analíticos da Política Anti-Inflacionária \thanks{Texto Traduzido livremente por Luiz Mario Andrade com o auxílio da ferramenta Chat GPT 5.4}}
\author{Paul A. Samuelson e Robert M. Solow}
\date{Maio de 1960}

\begin{document}

\maketitle

\begin{center}
    \textbf{PROBLEMA DE ALCANÇAR E MANTER UM NÍVEL ESTÁVEL DE PREÇOS}
\end{center}

\section*{I}

Assim como se diz que os generais estão sempre lutando a guerra errada, os economistas têm sido acusados de combater a inflação errada. Assim, na época da alta dos preços americanos de 1946-48, muita atenção foi focada nas sucessivas rodadas de aumentos salariais resultantes da negociação coletiva. No entanto, provavelmente a maioria dos economistas concorda agora que esse primeiro aumento de preços no pós-guerra foi atribuível principalmente à pressão da demanda (\textit{demand-pull}) resultante das acumulações de ativos líquidos e necessidades adiadas durante o período de guerra.

Essa ênfase na demanda foi de certa forma reforçada pela disparada de preços da Guerra da Coreia após meados de 1950. Mas, no momento em que a pressão de custos (\textit{cost-push}) estava se tornando desacreditada como teoria da inflação, deparamo-nos com o fenômeno bastante intrigante da elevação gradual (\textit{upward creep}) dos preços de 1955-58, que pareceu ocorrer na parte final do período, apesar da crescente capacidade ociosa, mercados de trabalho frouxos, crescimento real lento e nenhuma grande flutuabilidade aparente na demanda agregada.

Não é de admirar, então, que os economistas venham debatendo as possíveis causalidades envolvidas na inflação: demanda (\textit{demand-pull}) versus custos (\textit{cost-push}); pressão salarial (\textit{wage-push}) versus a mais geral "inflação de vendedores" de Lerner; e a nova teoria de Charles Schultze de inflação por "deslocamento de demanda" (\textit{demand-shift}). Propomo-nos a fazer um breve levantamento das questões. Em vez de nos pronunciarmos sobre a questão terrivelmente difícil de qual é exatamente o melhor modelo a usar para explicar o passado recente e prever o futuro provável, tentaremos enfatizar os tipos de evidência que podem ajudar a decidir entre as teorias conflitantes. E nos ocuparemos de algumas implicações de política que surgem das diferentes hipóteses analíticas.

\textit{História do Debate: A Teoria Quantitativa e a Demanda}. Os economistas pré-clássicos cresceram em um ambiente de preços secularmente ascendentes. E mesmo antes de Adam Smith, havia se desenvolvido a crença em pelo menos uma teoria quantitativa simplificada. Mas foi no pensamento neoclássico de Walras, Marshall, Fisher e outros que esta versão especial da determinação pela demanda do nível absoluto de preços monetários e custos atingiu sua forma mais desenvolvida.

Podemos simplificar excessivamente a doutrina da seguinte forma: Os produtos reais, insumos e preços relativos de bens e fatores podem ser pensados como determinados por um conjunto de equações competitivas que são independentes do nível absoluto de preços. Como em um sistema de escambo, o nível absoluto de todos os preços é indeterminado e inessencial devido às propriedades de "homogeneidade relativa" dessas relações de mercado. Para fixar o fator de escala absoluto, podemos, se quisermos, introduzir uma moeda neutra. Tal moeda, ao contrário do café ou do sabão, sendo valorizada apenas pelo que comprará e não por sua utilidade intrínseca, terá sua demanda exatamente dobrada se houver uma duplicação exata de todos os preços. Devido a essa importante "homogeneidade de escala", fixar o total dessa moeda irá, quando aplicado ao nosso sistema real já determinado de produtos, fatores e preços relativos, fixar o nível absoluto de todos os preços; e mudanças no total dessa moeda devem necessariamente corresponder a novos equilíbrios de preços absolutos que se moveram em proporção exata, com os preços relativos e todas as magnitudes reais permanecendo inalterados.\footnote{Mas, como Hume cedo reconheceu, os períodos de preços crescentes pareciam dar um estímulo transitório à economia, à medida que os buscadores de lucro ativos ganhavam vantagem à custa dos setores de renda fixa, credores e assalariados mais inertes. O outro lado desta tese de Hume é talvez exemplificado pelo fato de que as décadas de deflação pós-Guerra Civil foram também períodos de forte agitação social e de booms relativamente fracos e longos períodos de depressões mais pesadas do que a média — como estudos anteriores do National Bureau sugeriram.}

Como Patinkin e outros mostraram, as doutrinas acima são bastante simplificadas, pois não analisam totalmente as complexidades envolvidas na demanda por moeda; em vez disso, ignoram mudanças importantes (e previsíveis) em coeficientes de proporcionalidade como a velocidade de circulação. Mas, na Primeira Guerra Mundial, essa versão particular e estreita da inflação por demanda havia mais ou menos triunfado. A alta de preços no tempo de guerra era geralmente analisada em termos de aumentos na oferta monetária total. E a inflação alemã do pós-guerra foi compreendida por economistas não alemães em termos semelhantes.

Mas nem todos os economistas concordam com tudo. Assim como Tooke explicou de forma eclética o aumento de preços napoleônico parcialmente em termos do aumento de impostos, fretes e outros custos induzidos pela guerra, Harold G. Moulton e outros optaram por atribuir os aumentos de preços da Primeira Guerra Mundial a aumentos prévios no custo de produção. E não é sem significado que o grande neoclássico Wicksell expressou, nos últimos anos de sua vida, algumas dúvidas sobre a versão usual dos movimentos de preços em tempo de guerra, dando grande ênfase aos movimentos na velocidade da moeda induzidos pela escassez de bens em tempo de guerra.

É claro que os escritores neoclássicos não teriam negado a igualdade necessária entre custos competitivos e preços. Mas eles teriam considerado superficial tomar o nível de custos monetários como uma variável predeterminada. Em vez disso, argumentariam que preços e custos de fatores são simultaneamente determináveis em mercados competitivos interdependentes; e se o nível da oferta monetária total fosse mantido suficientemente sob controle, então o nível de preços poderia ser estabilizado, com quaisquer aumentos nos custos reais ou quaisquer diminuições na produção sendo compensados por pressão retroativa suficiente sobre os preços dos fatores, de modo a deixar os custos monetários finais e os preços, em média, inalterados.

Muitos escritores foram erroneamente além do argumento acima para conclusões insustentáveis, como a seguinte: um aumento nas despesas de defesa compensado por, digamos, impostos sobre o consumo não pode elevar o nível de preços se a quantidade de moeda for mantida constante; em vez disso, deve resultar em uma diminuição suficiente nos salários e outros custos de fatores para compensar exatamente o aumento nos custos tributários. Na verdade, porém, tal mudança na política fiscal poderia ser interpretada como uma redução na parcimônia combinada pública e privada; com $M$ constante, tenderia a aumentar o volume de gastos totais, exercendo pressão ascendente sobre as taxas de juros e induzindo um aumento na velocidade da moeda, resultando provavelmente em um nível de equilíbrio de preços mais elevado. Para fazer os preços voltarem ao nível anterior, seria necessária, mesmo dentro da estrutura de um modelo neoclássico estritamente competitivo, uma redução determinada na oferta monetária prévia. (Isso ilustra o perigo de passar da hipótese inocente — de que uma mudança equilibrada em todos os preços pode, a longo prazo, ser consistente com nenhuma mudança substantiva nas relações reais — para uma interpretação excessivamente simples de uma mudança complicada que está realmente ocorrendo na realidade histórica.)

Embora o exemplo acima de um aumento de preço induzido por impostos que ocorre dentro de um modelo neoclássico estrito possa ser denominado um caso de pressão de custos (\textit{cost-push}) em vez de demanda (\textit{demand-pull}), ele não representa realmente o mesmo fenômeno que encontraremos em nossa discussão posterior sobre pressão de custos. Isso pode ser visto mais facilmente pela observação de que, se alguém insistisse em manter os preços estáveis, os métodos convencionais de redução da demanda funcionariam muito bem, dentro do modelo neoclássico, para compensar tal pressão de custos.

\textit{Demanda à la Keynes}. Além da teoria quantitativa neoclássica, há uma segunda versão de inflação de demanda associada às teorias de Keynes. Antes e durante a Grande Depressão, os economistas ficaram impressionados com as fricções e rigidezes institucionais que causavam inflexibilidades descendentes em salários e preços e que tornavam tais movimentos deflacionários socialmente dolorosos. A \textit{Teoria Geral} de Keynes pode, se estivermos dispostos a simplificar excessivamente, ser pensada como um modelo sistemático que utiliza a inflexibilidade descendente de salários e preços para converter qualquer redução nos gastos monetários em uma redução real na produção e no emprego, em vez de uma redução equilibrada em todos os preços e custos de fatores. (Isso é excessivamente simples por pelo menos as seguintes razões: na versão pessimista de depressão de alguns keynesianos, uma hiperdeflação de salários e preços não teria tido efeitos substantivos na restauração do emprego e da produção, devido à elasticidade infinita da preferência pela liquidez e/ou elasticidade zero da demanda de investimento; na forma geral da \textit{Teoria Geral}, e particularmente após os efeitos Pigou do valor real da moeda terem sido incorporados, se fosse possível engendrar uma redução massiva de salários e custos, teria havido alguns efeitos estimulantes sobre o consumo, o investimento e a produção real; finalmente, uma teoria neoclássica cuidadosa, que levasse devidamente em conta as rigidezes e que analisasse as mudanças induzidas na velocidade de forma sofisticada, poderia também ter emergido com conclusões válidas semelhantes.)

Embora as teorias keynesianas possam ser consideradas diferentes das teorias neoclássicas com respeito à análise da deflação, o próprio Keynes estava disposto a assumir que o alcance do pleno emprego tornaria os preços e salários flexíveis para cima. Em \textit{Como Pagar pela Guerra} (1939), ele desenvolveu uma teoria da inflação que era muito parecida com a teoria neoclássica em sua ênfase na pressão da demanda do gasto agregado, embora diferisse daquela teoria em sua ênfase no fluxo total de gastos em vez de no estoque de moeda. Sua teoria da "inflação dos demandantes" derivava principalmente do fato de que o governo, mais os investidores, mais os consumidores desejam, em termos reais entre eles, mais de 100 por cento da produção produzível disponível em tempo de guerra ou de boom. Assim, os preços têm que subir para "enganar" os que gastam lentamente em relação às suas parcelas desejadas. Mas o aumento de preços fecha o hiato inflacionário apenas temporariamente, pois o nível de preços mais alto gera rendas mais altas para todos e o hiato real reabre-se continuamente. E assim a inflação prossegue, a uma taxa determinada pelo grau de deslocamento para os lucros, pela rapidez e extensão dos ajustes salariais ao custo de vida crescente e, finalmente, pela medida em que as receitas fiscais progressivas aumentam o suficiente para fechar o hiato. E, podemos acrescentar, a firmeza do banco central em limitar a oferta monetária pode, em última análise, aumentar tanto o aperto do crédito e reduzir tanto os saldos reais a ponto de equilibrar os gastos de consumo e investimento com os recursos civis disponíveis em algum patamar de preços mais elevado.

\textit{Teorias de Inflação por Pressão de Custos e Deslocamento de Demanda}. Em sua forma mais rígida, o modelo neoclássico exigiria que os salários caíssem sempre que houvesse desemprego de mão de obra e que os preços caíssem sempre que existisse capacidade ociosa, no sentido de que o custo marginal da produção que as empresas vendem é menor do que os preços que recebem. Um modelo mais eclético de competição imperfeita nos mercados de fatores e produtos é necessário para explicar o fato de aumentos de preços e salários antes que o pleno emprego e a plena capacidade tenham sido alcançados.

Da mesma forma, o modelo de Keynes, que assume rigidez salarial mesmo diante de um equilíbrio de subemprego, baseia-se em várias suposições de competição imperfeita. E quando reconhecemos que, consideravelmente antes de o pleno emprego de mão de obra e de plantas ser alcançado, os preços e salários modernos parecem mostrar uma tendência a derivar para cima irreversivelmente, vemos que o sistema keynesiano simples deve ser modificado ainda mais na direção de um modelo de competição imperfeita.

Ora, o fato de um modelo econômico envolver em algum grau competição imperfeita não implica necessariamente que os conceitos de mercados competitivos deem pouca visão sobre o comportamento de preços relativos, alocações de recursos e lucratividades. Para algum grau de aproximação, o modelo competitivo pode lançar luz sobre essas importantes magnitudes reais e, para este propósito, poderíamos nos contentar em usar o modelo competitivo. Mas para explicar a possível inflação de pressão de custos, pareceria mais econômico, desde o início, reconhecer que a competição imperfeita é a essência do problema e abandonar as suposições de competição perfeita.

Uma vez feito isso, reconhecemos a possibilidade qualitativa da inflação por pressão de custos. Assim como salários e preços podem ser rígidos diante do desemprego e da sobrecapacidade, também podem estar sendo empurrados para cima além do que pode ser explicado em termos de níveis e deslocamentos na demanda. Mas em que medida esses elementos são importantes na explicação do comportamento dos preços de qualquer período torna-se uma questão quantitativa importante. Não se deve esperar sempre que, observando o comportamento de uma economia ao longo de um determinado período, sejamos capazes de fazer uma separação muito boa do seu aumento de preços em elementos de demanda e de custo. Simplesmente não podemos realizar os experimentos controlados necessários para fazer tal separação; e a Mãe Natureza pode não ter nos dado economicamente a dispersão e variação necessárias como substituto para experimentos controlados, se quisermos fazer uma identificação aproximada das forças casuais em jogo.

Muitos economistas argumentaram que a pressão de custos foi importante no próspero período de 1951-53, mas que seus efeitos sobre os preços médios foram mascarados pela queda nos preços flexíveis das matérias-primas. Mas novamente em 1955-58, ela se manifestou apesar do fato de que, em boa parte deste período, parecia haver pouca evidência de pleno emprego generalizado e excesso de demanda. Alguns defensores desta visão atribuem a pressão a aumentos salariais engendrados unilateralmente por sindicatos fortes. Outros dão tanto ou mais peso à ação cooperativa de todos os vendedores — trabalho organizado e não organizado, gestões semimonopsonistas, vendedores oligopolistas em mercados de produtos imperfeitos — que elevam preços e custos na tentativa de cada um manter ou elevar sua parcela na renda nacional e que, entre si, ao tentar obter mais de 100 por cento do produto disponível, criam a "inflação de vendedores".

Uma variante da pressão de custos é fornecida pela teoria da inflação por "deslocamento de demanda" de Charles Schultze. A força da demanda em certos setores da economia — por exemplo, indústrias de bens de capital em 1955-57 — eleva preços e salários neles. Mas em outros lugares, embora a demanda não seja particularmente forte, a inflexibilidade descendente impede que os preços caiam, e o poder de mercado pode até engendrar um movimento de preço-salário imitativo, em algum grau, dos setores com forte demanda. O resultado é uma deriva ascendente nos preços médios — com a sugestão de que as políticas monetária e fiscal restritivas o suficiente para evitar um aumento médio de preços teriam de ser tão restritivas a ponto de produzir um nível considerável de desemprego e uma queda significativa na produção.

\section*{II}

\textit{Verdades e Consequências: O Problema da Identificação}. As teorias concorrentes (embora imperfeitamente concorrentes) da inflação parecem ser hipóteses genuinamente diferentes sobre fatos observáveis. Nesse caso, dever-se-ia ser capaz de distinguir empiricamente entre inflação de custo e de demanda. Quais são as marcas distintivas? Se eu acredito em pressão de custos, o que devo esperar encontrar nos fatos que não esperaria encontrar se fosse um crente na pressão de demanda? A última cláusula é importante. Não bastará apontar circunstâncias que acompanharão qualquer inflação, independentemente da causa. Um teste deve ter o que os estatísticos chamam de poder contra as principais hipóteses alternativas.

Por mais triviais que essas observações possam parecer, elas precisam ser feitas. Os clichês da discussão popular caem na armadilha repetidamente. Embora tenham sido pisoteados com frequência por especialistas, os erros revivem. Tomaremos o tempo para apontar o dedo mais uma vez. Fazemos isso porque queremos dar um passo adiante e argumentar que este problema de identificação é extremamente difícil. O que parece à primeira vista ser formas sutis e confiáveis de distinguir a inflação induzida por custos da induzida por demanda revela-se longe de ser infalível. Na verdade, somos levados a acreditar que os dados agregados, registrando os detalhes \textit{ex post} de transações concluídas, podem ser, na maioria das circunstâncias, bastante insuficientes. Pode ser necessário primeiro desagregar.

\textit{Falácias Comuns}. O erro mais simples — a ser encontrado em quase qualquer discussão de jornal sobre o assunto — é a crença de que, se os salários nominais sobem mais rápido que a produtividade, temos um sinal seguro de inflação de custos. É claro que a verdade é que, na forma mais pura de inflação por excesso de demanda, os salários subirão mais rápido que a produtividade; a única alternativa é que todo o aumento no valor de um produto fixo seja drenado para os lucros, sem que isso transborde para o mercado de trabalho para elevar os salários ainda mais. Este erro é às vezes misturado com a crença de que é possível, durante longos períodos, que indústrias com rápido aumento de produtividade paguem salários cada vez mais altos do que aquelas onde a produção por homem-hora cresce lentamente. Tal diferencial persistente e crescente provavelmente acabaria por alterar a composição de habilidades ou qualidades da força de trabalho nas diferentes indústrias, o que lança dúvidas sobre a comparação original de produtividade.

Às vezes veem-se afirmações no sentido de que aumentos na despesa mais rápidos que os aumentos na produção real significam necessariamente inflação de demanda. É aritmética simples que a despesa superando a produção por si só soletra apenas aumentos de preços e não fornece nenhuma evidência sobre a fonte ou causa da inflação. Muito do discurso sobre "muito dinheiro perseguindo poucos bens" é desse tipo.

Uma versão mais solene da falácia diz: Um aumento na despesa pode ocorrer apenas através de um aumento no estoque de moeda ou de um aumento na velocidade de circulação. Portanto, as únicas causas possíveis da inflação são $M$ e $V$ e não precisamos procurar mais.

\textit{Dificuldades Adicionais}. É mais desconcertante perceber que mesmo alguns dos testes empíricos sugeridos na literatura profissional podem ter pouco ou nenhum poder de corte para distinguir a inflação de custo da de demanda.

Pensa-se automaticamente em olhar para as relações de tempo. Os aumentos salariais parecem preceder os aumentos de preços? Então o aumento geral nos preços é causado pela pressão salarial. Os aumentos de preços parecem preceder os aumentos salariais? Então, mais provavelmente, a inflação é do tipo excesso de demanda, e os salários estão sendo puxados para cima por uma demanda vigorosa por mão de obra ou estão respondendo a aumentos anteriores no custo de vida. Existem pelo menos três dificuldades com este argumento. A primeira é sugerida pela substituição de "aumento salarial" por "galinha" e "aumento de preço" por "ovo". O problema é que não temos um padrão inicial normal a partir do qual medir, nenhum nível de preço que sempre existiu e ao qual todos se ajustaram; de modo que um aumento salarial, se ocorrer, deve ser autônomo e não uma resposta a alguma mudança prévia na demanda por mão de obra. Como ilustração da dificuldade de inferência, considere os ganhos horários médios na indústria siderúrgica básica. Eles subiram, em relação a toda a manufatura a partir de 1950, incluindo alguns períodos em que os mercados de trabalho não estavam apertados. Isso representou uma pressão salarial autônoma? Ou foi antes um ajuste retardado ao declínio dos salários do aço em relação a toda a manufatura, que ocorreu durante a guerra, presumivelmente como consequência da eficiência diferencial do controle salarial? E por que deveríamos tomar 1939 ou 1941 como padrão para salários relativos? E assim por diante.

Um problema relacionado é que, em uma economia estreitamente interdependente, os efeitos podem preceder as causas. Os preços podem começar a subir porque se espera que as taxas salariais subam. E, mais importante, à medida que os aumentos de salários e preços se propagam pela economia, a agregação pode facilmente distorcer as relações temporais aparentes.

Mas mesmo se pudéssemos encontrar a aparência de um experimento controlado, se após um período de estabilidade em ambos notássemos um aumento salarial para um novo patamar seguido de um aumento de preço, o que poderíamos concluir com segurança? Seria imensamente tentador fazer o diagnóstico óbvio de pressão salarial. Mas considere a seguinte cadeia hipotética de eventos: Os preços em mercados de produtos imperfeitos respondem apenas a mudanças nos custos. Os mercados de trabalho são perfeitamente competitivos na prática, e o salário monetário move-se rapidamente em resposta a deslocamentos na demanda por trabalho. Assim, qualquer explosão de excesso de demanda, gastos do governo, digamos, causaria um aumento na demanda por trabalho; os salários seriam puxados para cima; e só então os preços das mercadorias subiriam em resposta ao aumento do custo. Portanto, o diagnóstico óbvio pode estar errado. No meio, se fôssemos perspicazes, poderíamos notar um estreitamento temporário das margens e, com essa informação, poderíamos montar a história.

Considere outra inferência sofisticada. Em um mercado único, o preço pode subir ou porque a curva de demanda se desloca para a direita ou porque a curva de oferta se desloca para a esquerda em consequência de aumentos de custos. Mas no primeiro caso, a produção deve aumentar; no segundo caso, diminuir. Não poderíamos raciocinar, então, que se os preços sobem, setor por setor, com as produções, a pressão da demanda deve estar em ação? Muito provavelmente podemos, mas não com certeza. Em primeiro lugar, como Schultze argumentou, é possível que certos setores enfrentem excesso de demanda, sem que haja pressão agregada; esses setores mostrarão de fato fortes aumentos de preços e aumentos na produção (ou pressão sobre a capacidade). Mas, num sentido real, a fonte da inflação é a falha de outros setores, nos quais se desenvolve excesso de capacidade, de diminuir seus preços suficientemente. E isso pode ser uma consequência de "preços administrados", margens de lucro (\textit{markups}) rígidas, salários rígidos e toda a parafernália da "nova" inflação.

Para ir mais fundo, o raciocínio que estamos escrutinando pode falhar porque é ilegítimo, mesmo desta forma indústria por indústria, usar o raciocínio de equilíbrio parcial. Suponha que os salários subam. Somos levados a esperar uma diminuição na produção. Mas no mundo moderno, todos ou a maioria dos salários estão aumentando. Nem esta é a primeira vez que o fazem. E no passado, aumentos generalizados de salários e preços não resultaram em qualquer diminuição na demanda real agregada — talvez o contrário. Assim, mesmo em uma única indústria, as curvas de oferta e demanda podem não ser independentes. O deslocamento nos custos é acompanhado, de fato pode provocar, um deslocamento compensatório na curva de demanda subjetivamente vista pela indústria. E assim os preços podem subir sem declínio e possivelmente com um aumento na produção. Se houver algo nesta linha de pensamento, pode ser que uma das causas importantes da inflação seja — a inflação.

\textit{A Necessidade de Detalhe}. Nestes últimos parágrafos, estivemos argumentando contra a tentativa de diagnosticar a fonte da inflação a partir de agregados. Também sugerimos que às vezes os sintomas reveladores podem ser descobertos se não olharmos para os totais, mas para as partes. Essa sugestão ganha força quando reconhecemos, como devemos, que o mesmo aumento geral de preços pode facilmente ser consequência de diferentes causas em diferentes setores. Uma teoria monolítica pode ter sua simplicidade e estilo crivados de exceções. Existe alguma razão, além do desejo de simetria, para acreditarmos que o mesmo raciocínio deve explicar o aumento acima da média no preço dos serviços e o aumento acima da média no preço das máquinas desde 1951 ou desde 1949? Os preços dos serviços públicos indubitavelmente foram mantidos baixos durante a guerra pelo processo regulatório; e os serviços acompanham a demanda elástica em relação à renda, acompanhada por um aumento de produtividade registrado abaixo da média. Um aumento de preço mais rápido que a média equivale à mudança de preço relativo corretiva que se esperaria. O fator principal no caso das máquinas, de acordo com um estudo recente do Comitê Econômico Conjunto, parece ter sido uma explosão de excesso de demanda ocasionada pelo boom de investimentos de meados da década de cinquenta. E para dar ainda uma terceira variante, Eckstein e Fromm em outro estudo do Comitê Econômico Conjunto sugerem que o aumento acima da média nos salários dos trabalhadores siderúrgicos e nos preços dos produtos siderúrgicos ocorreu diante de um mercado de trabalho e de produtos um pouco menos apertado do que em máquinas. Eles atribuem isso a um exercício conjunto de poder de mercado pelo sindicato e pela indústria. Certa ou errada, é um erro tático teórico negar essa possibilidade alegando que ela não pode explicar a história dos preços em outros setores.

\textit{Algumas Coisas Que Seria Bom Saber}. Existem pelo menos duas questões clássicas que são relevantes para o nosso problema e sobre as quais surpreendentemente pouco trabalho foi feito: Uma é o comportamento da demanda real sob condições inflacionárias e a outra é o comportamento dos salários nominais em relação ao nível de emprego. Comentamos brevemente sobre estas duas questões porque nos parece haver alguma dúvida de que equações de comportamento reversíveis comuns possam ser encontradas, e esta mesma dificuldade aponta para uma questão importante que mencionamos anteriormente: que um período de alta demanda e preços crescentes molda atitudes, expectativas, e até instituições de tal forma a inclinar o futuro em favor de mais inflação. Ao contrário de alguns outros economistas, não tiramos a conclusão firme de que, a menos que se ponha um fim firme, a taxa de aumento de preços deve acelerar. Deixamos como uma questão aberta: Pode ser que a inflação rastejante leve apenas à inflação rastejante.

A maneira padrão para um hiato inflacionário se esgotar antes da hiperinflação é o próprio processo de inflação reduzir as demandas reais. Os mecanismos, alguns duvidosos, outros não, são bem conhecidos: o deslocamento para o lucro, efeitos de saldo real, progressividade tributária, aperto nas rendas fixas. Se os aumentos de preços e salários tiverem esse efeito, então uma inflação por pressão de custos na ausência de excesso de demanda inflige desemprego e excesso de capacidade ao sistema. A disposição de suportar a redução da demanda real é uma medida da imperfeição dos mercados que permite a pressão de custos. Mas suponha que as demandas reais não se comportem desta maneira? Suponha que um aumento de salários e preços não tenha efeito na demanda real, ou um efeito insignificante, ou mesmo um leve efeito positivo? Então não apenas a punição não se materializará, mas toda a distinção entre pressão de custos e pressão de demanda começa a evaporar. Mas isso é possível? Os teóricos quantitativos mais antigos certamente o teriam negado; mas o aumento da velocidade entre 1955 e 1957 teria surpreendido um teórico quantitativo mais antigo.

Não sabemos se a demanda real se comporta desta maneira ou não. Mas pensamos ser importante perceber que quanto mais o passado recente for dominado pela inflação, pelo alto emprego e pela crença de que ambos continuarão, mais provável é que o processo de inflação preserve ou mesmo aumente a demanda real, ou mais pesadamente as autoridades monetárias e fiscais tenham de incidir sobre a demanda no interesse da estabilização de preços. A consciência da renda real é uma força poderosa. A pressão sobre os saldos reais devido aos preços altos será parcialmente aliviada pela expectativa de preços crescentes, desde que as taxas de juros em um mercado de capitais imperfeito não consigam acompanhar o ritmo. As mesmas expectativas induzirão professores, aposentados e outros a tentar conceber instituições para proteger suas rendas reais da erosão pelos preços mais altos. Na medida em que tiverem sucesso, suas demandas reais permanecerão inalteradas. À medida que o medo do desemprego prolongado desaparece e a experiência do pleno emprego passado acumula economias, os assalariados também podem manter seus gastos reais; e as mesmas forças podem aumentar substancialmente a propensão marginal a gastar a partir dos lucros, incluindo lucros retidos. Se houver algo nesta linha de pensamento, o problema empírico da verificação pode ser muito difícil, porque muito da experiência do passado é irrelevante para a hipótese. Mas seria bom saber.

\textit{A Escala Fundamental de Phillips Relacionando Desemprego e Mudanças Salariais}. Considere também a questão da relação entre as mudanças nos salários nominais e o grau de desemprego. Temos o interessante artigo de A. W. Phillips sobre a história do Reino Unido desde a Guerra Civil (a nossa Guerra Civil!). Suas descobertas são notáveis, mesmo que se discorde de suas interpretações.

Em primeiro lugar, o período 1861-1913, durante o qual o movimento sindical foi bastante fraco, mostra uma relação bastante próxima entre a variação percentual nas taxas salariais e a fração da força de trabalho desempregada. Deve-se levar em conta as mudanças acentuadas no custo de vida induzidas pelos preços de importação e a expectativa normal de que os salários estarão subindo mais rápido quando uma taxa de desemprego de 5 por cento é atingida na fase ascendente do que quando é atingida na fase descendente. Em segundo lugar, com pequenas exceções, a mesma relação que serve para 1861-1913 também parece servir razoavelmente bem para 1913-48 e 1948-57. E finalmente Phillips conclui que o nível de salário nominal se estabilizaria com 5 por cento de desemprego; e a taxa de aumento dos salários nominais seria mantida na taxa de 2-3 por cento de aumento da produtividade com cerca de 2 $1/2$ por cento da força de trabalho desempregada.

Curiosamente, nenhum estudo comparavelmente cuidadoso foi feito para os EUA. A nota de Garbarino de 1950 é dificilmente uma análise em escala real, e o tratamento de Schultze em sua monografia de primeira classe do Comitê Conjunto é casual demais. Há alguma evidência de que os EUA diferem do Reino Unido em pelo menos dois aspectos. Se existe tal relação caracterizando o mercado de trabalho americano, ela pode ter se deslocado um pouco nos últimos cinquenta ou sessenta anos. Em segundo lugar, há uma sugestão de que neste país pode levar de 8 a 10 por cento de desemprego para estabilizar os salários nominais.

Mas levaria 8 a 10 por cento de desemprego para sempre para estabilizar o salário nominal? Não é este tipo de relação também dependente fortemente da experiência lembrada? Suspeitamos que este seja outro caminho pelo qual um passado caracterizado por preços crescentes, alto emprego e recessões curtas e suaves provavelmente gerará um viés inflacionário — tornando o salário nominal mais rígido para baixo, talvez até perversamente inclinado a subir durante recessões na esperança de que as coisas logo melhorem.

Pode não haver tal relação para este país. Se houver, por que não parece ter o mesmo grau de invariância a longo prazo que a curva de Phillips para o Reino Unido? Que fatos geográficos, econômicos, sociológicos explicam a diferença entre os dois países? Existe uma diferença na mobilidade do trabalho nos dois países? As diferentes tolerâncias para o desemprego refletem diferenças no nível de renda, organização sindical ou o quê? Que decisões políticas poderiam concebivelmente levar a uma diminuição na taxa crítica de desemprego na qual os salários começam a subir ou a subir rápido demais? Claramente, um estudo cuidadoso deste problema poderia render dividendos generosos.

\section*{III}

\textit{Um Olhar Mais Atento aos Dados Americanos}. Apesar de todas as suas deficiências, pensamos que o diagrama de dispersão acompanhante na Figura 1 é útil. Embora não forneça respostas, ele pelo menos faz perguntas interessantes. Plotamos as variações percentuais anuais dos ganhos horários médios na manufatura, incluindo suplementos (dados de Rees), contra a porcentagem média anual da força de trabalho desempregada.

O primeiro defeito a notar são as diferentes coberturas representadas nos dois eixos. Duesenberry argumentou que os aumentos salariais no pós-guerra na manufatura por um lado e no comércio, serviços, etc., por outro, podem ter explicações bastante diferentes: poder sindical na manufatura e simples excesso de demanda nos outros setores. Provavelmente é verdade que, se tivéssemos uma taxa de desemprego apenas para a manufatura, ela seria um pouco mais alta durante os anos do pós-guerra do que o número agregado mostrado. Mesmo que uma declaração qualitativa como esta fosse verdadeira para todo o período, o peso crescente dos serviços no total poderia ainda criar um viés. Outro defeito é o nosso uso de incrementos e médias anuais, quando um estudo em escala real teria de olhar cuidadosamente para as nuances do tempo.

Um primeiro olhar para a dispersão é desencorajador; há pontos por todo o lugar. Mas talvez se possa notar alguns efeitos sistemáticos. Em primeiro lugar, os anos de 1933 a 1941 parecem ser \textit{sui generis}: os salários nominais subiram ou não caíram diante de um desemprego massivo. Pode-se atribuir isso ao funcionamento do New Deal (o aumento salarial de 20 por cento de 1934 deve representar os códigos da NRA); ou alternativamente poder-se-ia argumentar que em 1933 muito do desemprego se tornara estrutural, isolado do mercado de trabalho em funcionamento, de modo que, na prática, o eixo vertical deveria ser movido para a direita. Isso deixaria algo mais parecido com o padrão normal.

Os primeiros anos da Primeira Guerra Mundial também se comportam de forma atípica, embora não tanto quanto 1933-39. Isso pode refletir aumentos no custo de vida, a rapidez do aumento da demanda, um aperto especial na manufatura, ou os três.

Mas o grosso das observações — o período entre a virada do século e a primeira guerra, a década entre o fim daquela guerra e a Grande Depressão, e os dez ou doze anos mais recentes — todos mostram um padrão bastante consistente. As taxas salariais tendem a subir quando o mercado de trabalho está apertado, e quanto mais apertado, mais rápido. O que é mais interessante é a forte sugestão de que a relação, tal como é, deslocou-se para cima ligeiramente, mas visivelmente, nos anos quarenta e cinquenta. Por um lado, a primeira década do século e os anos vinte parecem ajustar-se ao mesmo padrão. Os salários na manufatura parecem estabilizar-se absolutamente quando 4 ou 5 por cento da força de trabalho está desempregada; e aumentos salariais iguais ao aumento da produtividade de 2 a 3 por cento ao ano é o padrão normal com cerca de 3 por cento de desemprego. Isso não é tão terrivelmente diferente dos resultados de Phillips para o Reino Unido, embora a relação se mantenha lá com uma consistência maior. Comentamos sobre isso abaixo.

Por outro lado, de 1946 até o presente, o padrão é bastante consistente e consistentemente diferente do período anterior. A taxa de desemprego anual variou apenas estreitamente, de 2,5 por cento em 1953 a 6,2 por cento em 1958. Dentro desse intervalo, como seria de esperar, os salários subiram mais rápido quanto menor a taxa de desemprego. Mas julgar-se-ia agora que seriam necessários mais do que 8 por cento de desemprego para impedir o aumento dos salários nominais. E eles subiriam de 2 a 3 por cento ao ano com 5 ou 6 por cento da força de trabalho desempregada.

Seria precipitado concluir que a relação que estivemos discutindo representa uma curva de oferta reversível para o trabalho ao longo da qual uma curva de demanda agregada desliza. Se assim fosse, os movimentos ao longo da curva poderiam ser apelidados de pressão de demanda padrão, e deslocamentos da curva poderiam representar as mudanças institucionais nas quais se baseiam as teorias de pressão de custos. O deslocamento aparente em nossa curva de Phillips poderia ser atribuído por alguns economistas ao novo poder de mercado dos sindicatos. Outros poderiam estar mais inclinados a acreditar que a expectativa de pleno emprego contínuo, ou pelo menos alto emprego, é suficiente para explicar tanto o deslocamento na curva de oferta, se for isso, quanto a disposição dos empregadores (conscientes de que o que obtêm de uma força de trabalho depende em parte do seu moral e da sua rotatividade) de pagar aumentos salariais em períodos de demanda temporariamente frouxa.

Esta última consideração, no entanto, lança uma dúvida real sobre a fácil identificação da relação como meramente um fenômeno de oferta de trabalho. Existem duas partes em uma negociação salarial.

\textit{Comparação EUA e Reino Unido}. Uma comparação da posição americana com as descobertas de Phillips para o Reino Unido é interessante por si só e também como um possível guia para a política. Qualquer coisa que desloque a relação para baixo diminui o preço em desemprego que deve ser pago quando se segue uma política de conter a taxa de aumento de salários e preços por pressão na demanda agregada.

Uma possibilidade é que a liderança sindical seja mais "responsável" no Reino Unido; na verdade, a política de contenção salarial do pós-guerra parece visível nos dados de Phillips. Mas há outras interpretações. É claro que quanto mais fracionado e imperfeito for um mercado de trabalho, maior terá que ser o excesso de oferta total de trabalho antes que a taxa salarial média se torne estável e menos apertada será a relação em qualquer caso. Mesmo um toque de inflexibilidade descendente (e o sindicalismo e os salários administrados certamente significam pelo menos isso) tornará este efeito de imobilidade mais pronunciado. Pareceria plausível que a pura compactação geográfica da economia inglesa torne seu mercado de trabalho mais perfeito que o nosso neste sentido. Além disso, os britânicos têm perseguido uma política mais deliberada de realocação da indústria para absorver bolsões de desemprego estrutural.

Isso sugere que qualquer política governamental que aumente a mobilidade do trabalho (geográfica e industrial) ou melhore o fluxo de informações no mercado de trabalho terá efeitos anti-inflacionários, além de ser desejável por outras razões. Uma abordagem mais rápida, mas a longo prazo provavelmente menos eficiente, poderia ser o governo direcionar a distribuição regional de seus gastos mais deliberadamente em termos da existência de desemprego local e excesso de capacidade.

Os dados ingleses mostram uma relação claramente não linear (hiperbólica) entre as mudanças salariais e o desemprego, refletindo a muito discutida inflexibilidade descendente. Nossos números americanos não contradizem isso, embora não contem uma história tão clara quanto os ingleses. Na medida em que esta não linearidade existe, como Duesenberry observou, um dado nível médio de desemprego ao longo do ciclo será compatível com uma taxa mais lenta de aumento salarial (e presumivelmente aumento de preço) quanto menos amplas forem as oscilações cíclicas do topo ao fundo.

Uma implicação menos óbvia deste ponto de vista é que uma política deliberada de baixa pressão para estabilizar o nível de preços pode ter um certo aspecto de autodestruição. Fica claro pela experiência que a mobilidade inter-regional e inter-industrial do trabalho depende fortemente da atração de oportunidades de emprego em outros lugares, mais do que da pressão do desemprego local. De fato, a imperfeição do mercado de trabalho é aumentada, com as consequências que esboçamos.

\section*{IV}

Concluímos que não é possível, com base no raciocínio a \textit{priori}, rejeitar a hipótese da pressão de demanda ou da pressão de custos, ou as variantes desta última, como o deslocamento de demanda. Também argumentamos que as identificações empíricas necessárias para distinguir entre essas hipóteses podem ser quase impossíveis a partir da experiência de macro-dados que está disponível para nós; e que, embora o uso de micro-dados possa lançar luz adicional sobre o problema, mesmo aqui a identificação é repleta de dificuldades e ambiguidades.

No entanto, há uma área onde o interesse político e o desejo de compreensão científica por si só se unem. Se, por uma política deliberada, se engendrasse uma redução considerável na demanda ou se recusasse a permitir o aumento na demanda que seria necessário para preservar o alto emprego, ter-se-ia um experimento que poderia esperar distinguir entre a validade da teoria da pressão de demanda e da pressão de custos, conforme reformularíamos operacionalmente essas teorias. Se uma pequena redução da demanda fosse seguida por grandes moderações na marcha dos salários e outros custos, de modo que o custo social de um índice de preços estável resultasse ser muito pequeno em termos de emprego e produção de alto nível sacrificados, então a hipótese da pressão de demanda teria recebido sua confirmação mais importante. Por outro lado, se a repressão suave da demanda não contivesse os aumentos de custos e preços de forma alguma ou apenas suavemente, de modo que um desemprego considerável tivesse de ser engendrado antes que a deriva ascendente do nível de preços pudesse ser evitada, então a hipótese da pressão de custos teria recebido sua confirmação mais importante. Se o resultado desta experiência se revelasse entre esses casos extremos — como nós mesmos esperaríamos — então um elemento de validade teria de ser concedido a ambas as visões; e, por mais maçante que seja ter de abraçar teorias ecléticas, estudiosos que desejassem ser realistas teriam de se preparar para fazê-lo.

É claro que estivemos falando levianamente de um vasto experimento. Na verdade, tal operação estaria repleta de implicações para o bem-estar social. Naturalmente, como estão confiantes de que seria um sucesso, os crentes na pressão de demanda deveriam acolher tal experimento. Mas, igualmente natural, os crentes na pressão de custos estariam totalmente contra tal economia de baixa pressão engendrada, uma vez que estão igualmente convencidos de que será um fracasso desastroso envolvendo muita dor social desnecessária. (Uma terceira escola, que acredita na pressão de custos, mas pensa que ela pode ser curada ou minimizada pela depressão ortodoxa da demanda, pensa que nossa falha em realizar este experimento estaria repleta de mal social em virtude do fato de que esperam que uma elevação gradual nos preços se transforme em um trote e depois em um galope.)

Nossa própria visão agora terá se tornado evidente. Quando traduzimos o diagrama de Phillips mostrando o padrão americano de aumento salarial contra o grau de desemprego em um diagrama relacionado mostrando os diferentes níveis de desemprego que seriam "necessários" para cada grau de mudança no nível de preços, chegamos a palpites como os seguintes:

1. Para que os salários aumentem não mais do que os 2 $1/2$ por cento ao ano característicos do nosso crescimento de produtividade, a economia americana pareceria, com base na experiência do século XX e do pós-guerra, ter de passar por algo como 5 a 6 por cento da força de trabalho civil desempregada. Esse tanto de desemprego pareceria ser o custo da estabilidade de preços nos anos imediatamente à frente.

2. Para alcançar o objetivo do não perfeccionista de uma produção alta o suficiente para nos dar não mais do que 3 por cento de desemprego, o índice de preços poderia ter de subir até 4 a 5 por cento ao ano. Esse tanto de aumento de preços pareceria ser o custo necessário de alto emprego e produção nos anos imediatamente à frente.

Tudo isso é mostrado em nossa modificação do nível de preços da curva de Phillips, Figura 2. O ponto A, correspondente à estabilidade de preços, é visto envolver cerca de 5 $1/2$ por cento de desemprego; enquanto o ponto B, correspondente a 3 por cento de desemprego, é visto envolver um aumento de preços de cerca de 4 $1/2$ por cento ao ano. Esperamos que a queda de braço da política nos acabe colocando, nos próximos anos, em algum lugar entre esses pontos selecionados. Provavelmente teremos algum aumento de preços e algum excesso de desemprego.

Além do aviso usual de que estes são simplesmente os nossos melhores palpites, devemos dar outra cautela. Toda a nossa discussão foi formulada em termos de curto prazo, lidando com o que pode acontecer nos próximos anos. Seria errado, porém, pensar que o nosso menu da Figura 2, que relaciona o preço alcançável e o comportamento do desemprego, manterá a mesma forma a longo prazo. O que fizermos de forma política durante os próximos anos pode fazer com que ele se desloque de uma maneira definida.

Assim, é concebível que, após terem produzido uma economia de baixa pressão, os crentes na pressão de demanda fiquem desapontados no curto prazo; isto é, os preços podem continuar a subir mesmo que o desemprego seja considerável. No entanto, pode ser que a demanda de baixa pressão aja sobre as expectativas salariais e outras de modo a deslocar a curva para baixo no longo prazo — de modo que, ao longo de uma década, a economia possa desfrutar de maior emprego com estabilidade de preços do que nossa estimativa atual indicaria.

Mas também o oposto é concebível. Uma economia de baixa pressão pode acumular dentro de si ao longo dos anos quantidades cada vez maiores de desemprego estrutural (o inverso do que aconteceu de 1941 a 1953 como resultado da forte demanda de guerra e do pós-guerra). O resultado seria um deslocamento ascendente do nosso menu de escolha, com mais e mais desemprego sendo necessário apenas para manter os preços estáveis.

Como não temos evidências conclusivas ou sugestivas sobre essas questões conflitantes, não tentaremos dar um julgamento sobre elas. Em vez disso, aventuramos o lembrete de que, nos anos imediatamente à frente, o nível de crescimento alcançado estará altamente correlacionado com o grau de pleno emprego e produção de alta capacidade.

Mas e quanto ao longo prazo? Se a taxa anual de progresso técnico fosse aproximadamente a mesma em uma economia de baixa e de alta pressão, então a perda inicial de produção ao ir para o estado de baixa pressão nunca seria compensada; no entanto, em termos relativos, o hiato inicial não cresceria, mas permaneceria constante com o passar do tempo. Se uma economia de baixa pressão pudesse ter sucesso em melhorar a eficiência de nossos fatores produtivos, parte da perda de crescimento poderia ser gradualmente compensada e poderia em tempo suficiente ser até mais do que eliminada. Por outro lado, se tal economia produzisse guerra de classes e conflito social e deprimisse o nível de pesquisa e progresso técnico, a perda de crescimento seria composta no longo prazo.

Um aviso final é necessário. Não entramos aqui na importante questão de quais reformas institucionais viáveis poderiam ser introduzidas para diminuir o grau de desarmonia entre o pleno emprego e a estabilidade de preços. Estas poderiam, é claro, envolver questões de amplo alcance, como controles diretos de preços e salários, legislação antissindical e antitruste, e uma série de outras medidas esperançosamente concebidas para mover as curvas de Phillips americanas para baixo e para a esquerda.

\end{document}